% Custom CV — Elogen (Les Ulis) — spontaneous application, R&D PEM electrolysis — 2026-11-12
% Base: resumes/library/electrochem_fr/source.tex (+ P3 block from process_fr)
\documentclass[10pt]{article}
\usepackage[letter]{nermine-cv}
\begin{document}
\cvheader{Moumen Nermine}{Étudiante en dernière année du cycle ingénieur en Chimie Industrielle}{nermine.moumen@insat.ucar.tn}{+21627040414}{linkedin.com/in/nermine-moumen}

\section{FORMATION}
\entry{Institut National des Sciences Appliquées et de Technologie (INSAT)}{2022–Présent}
\summary{Étudiante en Chimie Industrielle, spécialité Procédés et Environnement, avec une formation solide en cinétique électrochimique, modélisation électrochimique sous COMSOL, ANSYS Fluent et MATLAB, et génie des procédés. Orientée vers la production d'hydrogène par électrolyse PEM et l'interprétation du comportement électrochimique des systèmes. À la recherche d'un PFE de 6 mois à partir de février 2027.}

\section{EXPÉRIENCE PROFESSIONNELLE}
\entry{SEGULA Technologies - Stage en Modélisation Électrochimique des Batteries}{06/2026–08/2026}
\begin{bul}
\item Modélisation d'une cellule \textbf{Li-ion} (NMC811/Graphite-SiOx) via le Single Particle Model (SPM) implémenté sous \textbf{MATLAB}.
\item Résolution des équations aux dérivées partielles par \textbf{FVM} avec intégration temporelle de Crank-Nicolson.
\item Optimisation du maillage et des coefficients de diffusion via une analyse de données, réduisant l'erreur \textbf{RMSE de 202,13 mV à 40,5 mV}.
\item Validation des résultats par une approche itérative et une analyse critique avec les données expérimentales et \textbf{COMSOL}.
\end{bul}
\entry{Control Union - Stage en Contrôle Qualité}{08/2025–09/2025}
\begin{bul}
\item Analyse de lubrifiants (viscosité et métaux lourds par \textbf{ICP-OES}).
\item Préparation et interprétation de rapports d'analyse pour vérifier la conformité aux normes \textbf{ASTM D5185} et \textbf{ISO 17025}.
\end{bul}
\entry{ETAP - Stage en Caractérisation Analytique d'hydrocarbures et en Traitement des Eaux}{07/2024–08/2024}
\begin{bul}
\item Étude des méthodes de traitement des eaux pétrolières par échange d'ions et par coagulation-floculation.
\item Formation aux techniques de géochimie organique, incluant la chromatographie phase gazeuse (GC, GC-MS).
\end{bul}

\section{PROJETS ACADÉMIQUES}
\project{Modélisation de Piles à Combustible PEM et Électrolyseurs PEM sous COMSOL et ANSYS Fluent}
\begin{bul}
\item Modélisation de piles à combustible (écoulement diphasique, \textbf{effets thermiques}) à courant imposé puis à potentiel imposé.
\item Modélisation d'un \textbf{électrolyseur PEM 3D} pour la \textbf{production d'hydrogène} : cinétique électrochimique, conservation de charge, identification de la \textbf{réaction limitante (OER)}.
\end{bul}
\project{Traitement Électrochimique d'un Rejet Industriel par Électro-floculation}
\begin{bul}
\item Élimination des cations toxiques (Cd$^{2+}$, Zn$^{2+}$) d'un effluent simulé par électro-coagulation/floculation avec des électrodes d'aluminium : taux d'abattement de \textbf{85 \% (Cd$^{2+}$) et 57 \% (Zn$^{2+}$) en 20 min}.
\item Quantification de la cinétique d'élimination par \textbf{polarographie impulsionnelle} (courbes courant-potentiel) et loi de Fick.
\end{bul}
\project{Analyse techno-environnementale de la production d'éthylène : vapocraquage vs déshydratation du bioéthanol}
\begin{bul}
\item Simulation comparative de deux voies de production sous \textbf{Aspen Plus V14}, avec établissement des \textbf{bilans matière et énergie} et des diagrammes de flux \textbf{(PFD)} du procédé, pour évaluer la faisabilité technique de chaque filière.
\item Analyse de \textbf{Cycle de Vie} cradle-to-gate sous \textbf{OpenLCA} pour comparer les \textbf{impacts environnementaux} des deux voies.
\end{bul}

\section{COMPÉTENCES TECHNIQUES}
\begin{bul}
\item \textbf{Techniques de Laboratoire :} Chromatographie (HPLC, GC, CCM), Spectrométrie (MS, ICP-OES), Spectroscopie (RMN, FTIR, UV-Vis), Diffraction (XRD), Microscopie (MEB, MET), Électrochimie (Voltamétrie cyclique, polarographie).
\item \textbf{Simulation :} Aspen Plus et HYSYS, COMSOL Multiphysics\sep\textbf{Programmation :} \textit{MATLAB, Mathematica, Python}
\item \textbf{Données \& Bureautique :} Pack Microsoft Office, OpenLCA\sep\textbf{Normes :} ISO 9001, ISO 14001, ISO 45001, ISO 17025, ISO 14040
\end{bul}

\section{ACTIVITÉS EXTRASCOLAIRES}
\textbf{IEEE INSAT} — Chef d'organisation de TRSYP 3.0 ; Responsable des activités sociales – MENASYP ; Présidente de l'EMBS.\\
\textbf{THEATRO INSAT} — Membre de l'équipe de danse.\par

\section{LANGUES ET CERTIFICATS}
\twocol{Langues}{Anglais (Avancé), Français (Bilingue), Arabe (Natif).}{Certificats}{Aspen HYSYS – MATLAB}
\end{document}
